\documentclass[10pt,a4paper]{amsart}
\usepackage[margin=19mm]{geometry}
\usepackage{amsmath,amssymb,mathtools}
\usepackage[scr=rsfs,cal=euler]{mathalfa}
\usepackage{xfrac}
\usepackage[unicode,hidelinks,bookmarks=false]{hyperref}
\newcommand{\ie}{i.e.\ }
\newcommand{\cf}{cf.\ }
\newcommand{\resp}{respectively\ }
\newcommand{\Subsection}{\subsection}
\providecommand{\nobreakdash}{\nobreak\hbox{-}}
\setlength{\emergencystretch}{2em}
\multlinegap0pt
\usepackage{amssymb}
\usepackage[all]{xy}
\usepackage{mathtools}
\newtheorem{lemma}{Lemma}
\newcommand{\ep}{\varepsilon}
\newcommand{\eq}[1]{\begin{equation}#1\end{equation}}
\newcommand{\genlegendre}[4]{%
  \genfrac{(}{)}{}{#1}{#3}{#4}%
  \if\relax\detokenize{#2}\relax\else_{\!#2}\fi
}
\newcommand{\legendre}[3][]{\genlegendre{}{#1}{#2}{#3}}
\title{A note on Shintani's invariant}
\author{Bora Yalkinoglu}
\date{Source: arXiv:2408.07309v3 (CC BY 4.0)}
\begin{document}
\maketitle

\noindent\emph{Source and licence.} A note on Shintani's invariant. Version arXiv:2408.07309v3 (CC BY 4.0), distributed by arXiv under the Creative Commons Attribution 4.0 International licence, http://creativecommons.org/licenses/by/4.0/, as recorded in the arXiv licence metadata for this exact version and shown on the abstract page https://arxiv.org/abs/2408.07309v3. Full licence notice: \texttt{licenses/arxiv-2408.07309-CC-BY-4.0.txt}.\par\smallskip

\noindent\textit{Editorial scope.} Counted unit: the article's lemma epsilonf (source0.tex lines 337--340). The article's complete original proof of it is delivered with it (source0.tex lines 341--345, one \texttt{proof} environment of 5 lines). No mathematics is written, repaired or re-derived: the statement and the proof are the article's own lines, in the article's own notation, and no retained line is substituted or re-flowed.  No further source-local context is needed: every cross-reference of every retained line resolves inside the two delivered spans.  Nothing was omitted from any retained span, so the package carries no omission record.  No retained line contains a citation, so the wrapper carries no bibliography and no bibliography entry.  No retained line contains a figure environment, an image-inclusion command or drawing code, so no image is delivered and the package declares no asset dependency.  Editorial formatting only, in addition: an AMS \texttt{amsart} preamble that restates the article's own declarations and macros used by the retained lines, namely the environments \texttt{lemma} on separate counters, together with the standard packages those lines need.  The retained environments print in this document as Lemma 1 in order of appearance, whereas the article prints the same items with the numbering of its own full document.  The complete native source of the article as distributed in the arXiv e-print for this exact version is bundled unchanged as \texttt{sources/163764/source0.tex}.
\begin{lemma}
\label{epsilonf}
Let $p \in \mathbb N$ be a prime number not dividing the discriminant $\Delta$ of $K$. Then, with the usual notation for the Jacobi symbol, we have: \eq{g(p) \ \vert \ p-\legendre{\Delta}{p}.}
\end{lemma}

\begin{proof}
Recall $\ep = \tfrac{a + b\sqrt{d} } 2$, with $a,b \in \mathbb N$, and set $l = p-\legendre{\Delta}{p}$. When we write $e\equiv f$, we always mean $e \equiv f \text{ mod $p$}$. \\ \\
Let us start with the case $(p)$ is non-split, i.e., in particular $\legendre{\Delta}{p} = -1$. We need to show that \eq{p \ | \ 2^{p+1}(\ep^{p+1}-1).} As \eq{\binom{p}{k} \equiv \left\{ \begin{array}{c} 0 \text{ , if $0 < k < p$} \\  \  1 \text{ , if $k \in \{0,p\}$}  \end{array}  \right. ,} we have \eq{2^{p+1} \ep^{p+1} = (a+b\sqrt{d} )^p(a+b\sqrt{d} ) = a^{p+1}+b^{p+1} d^{(p+1)/2} +(a^p b + a b^p d^{(p-1)/2} ) \sqrt{d}.} From the assumptions, we know \eq{d^{(p+1)/2} \equiv -d \text{ and } d^{(p-1)/2} \equiv -1 .} Moreover, using Fermat's little theorem, we see \eq{a^{p+1}+b^{p+1} d^{(p+1)/2} \equiv a^2 -d b^2 \equiv 4} and \eq{a^p b + a b^p d^{(p-1)/2} \equiv ab - ab = 0.} This gives \eq{2^{p+1}\ep^{p+1} \equiv 4 } and thus $p \ \vert \ 2^{p+1}(\ep^{p+1}-1)$. \\ \\
Now we look at the case where $(p)$ splits. In this case we need to show that \eq{p \ \vert \ (\ep^{l}-1) (\ep^{-l}-1) = 2 - (\ep^{l} + \ep^{-l}) = 2- T_l(a).} We calculate \eq{2^l \, T_l(a) = 2 \sum_{k=0}^{l/2} \binom{l}{2k}a^{2k} (b^2d)^{l/2-k}.} If $\legendre{\Delta}{p} = 1$, there exists $1\leq r \leq l/2$ such that $r^2 \equiv d$. Using $\binom{p-1}{2k} \equiv 1 $ for $0 \leq k \leq l/2$ and Fermat's little theorem, we obtain \eq{2^l \, T_l(a) \equiv T_l(a) \equiv 2 (br)^l \sum_{k=0}^{l / 2}(a^2b^{-2}d^{-1})^k \equiv 2 \tfrac{(a^2b^{-2}d^{-1})^{l+1}-1}{(a^2b^{-2}d^{-1})-1} \equiv 2 \tfrac{(a^2b^{-2}d^{-1})-1}{(a^2b^{-2}d^{-1})-1} = 2.} If $\legendre{\Delta}{p}=-1$, we have again $d^{l/2}\equiv -d $. From \eq{\binom{p+1}{k} \equiv \left\{ \begin{array}{c} 0 \text{ ,if $2 \leq k \leq p-1$} \\  \  1 \text{, if $k \in \{0,1,p,p+1\}$}  \end{array}  \right. ,} together with Fermat's little theorem, we get \eq{2^l \, T_l(a) \equiv 2(a^l+b^l d^{l/2}) \equiv 2(a^2 - d b^2 ) = 8 .} Thus, we get $p \ \vert \ T_l(a) - 2$, finishing the proof.
\end{proof}

\end{document}
